% GENERATED FILE. Edit canonical lessons, then run npm run generate:books.
% canonical-source-hash: fnv1a64:4f65e84cb3c742f6
% canonical-lessons: ML-C287-palaharam, ML-C287-palampori, ML-C287-pappaya, ML-C287-orancu, ML-C287-appil

\chapter{Snack, Banana fritter, Papaya, Orange, Apple}
\label{ch:ml-snack-banana-fritter-papaya-orange-apple}
\hlchaptermodality{287}

\begin{chapteropening}
\textbf{By the end of this chapter:} \emph{I can say a snack, a banana fritter, a papaya, an orange and an apple.}

Complete the last lesson of chapter 287: I can say a snack, a banana fritter, a papaya, an orange and an apple.
\end{chapteropening}

\section[\texorpdfstring{\ml{പലഹാരം} (palahāraṁ)}{പലഹാരം (palahāraṁ)}]{\ml{പലഹാരം} (palahāraṁ) — a snack}

\label{lesson:ML-C287-palaharam}



\begin{quote}
Before the new one: say the Malayalam for a purse, then the Malayalam for a pocket.
\end{quote}

\subsection*{You'll want to know: \ml{പലഹാരം}}
\textbf{\ml{പലഹാരം}} — \emph{palahāraṁ} — \textquotedblleft{}a snack\textquotedblright{}.

\begin{cousinweb}
\textbf{\ml{പലഹാരം}} (\emph{palahāraṁ}) $\leftarrow$ Sanskrit \emph{phalāhāra}, \textquotedblleft{}eating fruit\textquotedblright{}, reshaped as a tadbhava, with Malayalam's noun ending \emph{-ṁ} (CDIAL 9060). Its parts are \emph{phala}, \textquotedblleft{}fruit\textquotedblright{}, and \emph{āhāra}, \textquotedblleft{}food\textquotedblright{}; Malayalam softened the \emph{ph} and widened the sense to any light snack.
\end{cousinweb}

\begin{grammarlens}[title={What you've built}]
One more word for this chapter.
\end{grammarlens}

\subsection*{Your turn}
\begin{itemize}
\raggedright
  \item \emph{Say it:} \emph{palahāraṁ}
  \item \emph{Say it:} \emph{palahāraṁ}, once more
  \item \emph{Say it:} say \emph{palahāraṁ}
  \item \emph{Recall:} say the Malayalam for a purse, then the Malayalam for a pocket, then say \emph{palahāraṁ} again
\end{itemize}

\begin{tcolorbox}[breakable,colback=teal!4,colframe=teal!35!black,title={Before you move on}]
What does \ml{പലഹാരം} mean? (\textquotedblleft{}A snack\textquotedblright{}.) Say it once more.
\end{tcolorbox}
\section[\texorpdfstring{\ml{പഴംപൊരി} (paḻampori)}{പഴംപൊരി (paḻampori)}]{\ml{പഴംപൊരി} (paḻampori) — a banana fritter}

\label{lesson:ML-C287-palampori}



\begin{quote}
Before the new one: say the Malayalam for a pocket, then the Malayalam for a snack.
\end{quote}

\subsection*{You'll want to know: \ml{പഴംപൊരി}}
\textbf{\ml{പഴംപൊരി}} — \emph{paḻampori} — \textquotedblleft{}a banana fritter\textquotedblright{}.

\begin{cousinweb}
\textbf{\ml{പഴംപൊരി}} (\emph{paḻampori}) is native Dravidian, \textquotedblleft{}fried fruit\textquotedblright{}: \emph{paḻam}, \textquotedblleft{}ripe fruit\textquotedblright{} (DEDR 4004), and \emph{pori}, from the verb \textquotedblleft{}to fry, parch\textquotedblright{} (DEDR 4537, with Tamil \emph{pori}). DEDR 4004 compares Sanskrit \emph{phala}.
\end{cousinweb}

\begin{grammarlens}[title={What you've built}]
One more word for this chapter.
\end{grammarlens}

\subsection*{Your turn}
\begin{itemize}
\raggedright
  \item \emph{Say it:} \emph{paḻampori}
  \item \emph{Say it:} \emph{paḻampori}, once more
  \item \emph{Say it:} say \emph{paḻampori}
  \item \emph{Recall:} say the Malayalam for a pocket, then the Malayalam for a snack, then say \emph{paḻampori} again
\end{itemize}

\begin{tcolorbox}[breakable,colback=teal!4,colframe=teal!35!black,title={Before you move on}]
What does \ml{പഴംപൊരി} mean? (\textquotedblleft{}A banana fritter\textquotedblright{}.) Say it once more.
\end{tcolorbox}
\section[\texorpdfstring{\ml{പപ്പായ} (pappāya)}{പപ്പായ (pappāya)}]{\ml{പപ്പായ} (pappāya) — a papaya}

\label{lesson:ML-C287-pappaya}



\begin{quote}
Before the new one: say the Malayalam for a snack, then the Malayalam for a banana fritter.
\end{quote}

\subsection*{You'll want to know: \ml{പപ്പായ}}
\textbf{\ml{പപ്പായ}} — \emph{pappāya} — \textquotedblleft{}a papaya\textquotedblright{}.

\begin{cousinweb}
\textbf{\ml{പപ്പായ}} (\emph{pappāya}) is \emph{papaya}, a word from a language of the Caribbean that travelled with the fruit through Spanish and Portuguese.
\end{cousinweb}

\begin{grammarlens}[title={What you've built}]
One more word for this chapter.
\end{grammarlens}

\subsection*{Your turn}
\begin{itemize}
\raggedright
  \item \emph{Say it:} \emph{pappāya}
  \item \emph{Say it:} \emph{pappāya}, once more
  \item \emph{Say it:} say \emph{pappāya}
  \item \emph{Recall:} say the Malayalam for a snack, then the Malayalam for a banana fritter, then say \emph{pappāya} again
\end{itemize}

\begin{tcolorbox}[breakable,colback=teal!4,colframe=teal!35!black,title={Before you move on}]
What does \ml{പപ്പായ} mean? (\textquotedblleft{}A papaya\textquotedblright{}.) Say it once more.
\end{tcolorbox}
\section[\texorpdfstring{\ml{ഓറഞ്ച്} (ōṟañcŭ)}{ഓറഞ്ച് (ōṟañcŭ)}]{\ml{ഓറഞ്ച്} (ōṟañcŭ) — an orange}

\label{lesson:ML-C287-orancu}



\begin{quote}
Before the new one: say the Malayalam for a banana fritter, then the Malayalam for a papaya.
\end{quote}

\subsection*{You'll want to know: \ml{ഓറഞ്ച്}}
\textbf{\ml{ഓറഞ്ച്}} — \emph{ōṟañcŭ} — \textquotedblleft{}an orange\textquotedblright{}.

\begin{cousinweb}
\textbf{\ml{ഓറഞ്ച്}} (\emph{ōṟañcŭ}) $\leftarrow$ English \emph{orange}, which goes back through Arabic and Persian to the Sanskrit word for the orange tree.
\end{cousinweb}

\begin{grammarlens}[title={What you've built}]
One more word for this chapter.
\end{grammarlens}

\subsection*{Your turn}
\begin{itemize}
\raggedright
  \item \emph{Say it:} \emph{ōṟañcŭ}
  \item \emph{Say it:} \emph{ōṟañcŭ}, once more
  \item \emph{Say it:} say \emph{ōṟañcŭ}
  \item \emph{Recall:} say the Malayalam for a banana fritter, then the Malayalam for a papaya, then say \emph{ōṟañcŭ} again
\end{itemize}

\begin{tcolorbox}[breakable,colback=teal!4,colframe=teal!35!black,title={Before you move on}]
What does \ml{ഓറഞ്ച്} mean? (\textquotedblleft{}An orange\textquotedblright{}.) Say it once more.
\end{tcolorbox}
\section[\texorpdfstring{\ml{ആപ്പിൾ} (āppiḷ)}{ആപ്പിൾ (āppiḷ)}]{\ml{ആപ്പിൾ} (āppiḷ) — an apple}

\label{lesson:ML-C287-appil}



\begin{quote}
Before the new one: say the Malayalam for a papaya, then the Malayalam for an orange.
\end{quote}

\subsection*{You'll want to know: \ml{ആപ്പിൾ}}
\textbf{\ml{ആപ്പിൾ}} — \emph{āppiḷ} — \textquotedblleft{}an apple\textquotedblright{}.

\begin{cousinweb}
\textbf{\ml{ആപ്പിൾ}} (\emph{āppiḷ}) $\leftarrow$ English \emph{apple}.
\end{cousinweb}

\begin{grammarlens}[title={What you've built}]
One more word for this chapter.
\end{grammarlens}

\subsection*{Your turn}
\begin{itemize}
\raggedright
  \item \emph{Say it:} \emph{āppiḷ}
  \item \emph{Say it:} \emph{āppiḷ}, once more
  \item \emph{Say it:} say \emph{āppiḷ}
  \item \emph{Recall:} say the Malayalam for a papaya, then the Malayalam for an orange, then say \emph{āppiḷ} again
\end{itemize}

\begin{tcolorbox}[breakable,colback=teal!4,colframe=teal!35!black,title={Before you move on}]
What does \ml{ആപ്പിൾ} mean? (\textquotedblleft{}An apple\textquotedblright{}.) Say it once more.
\end{tcolorbox}
